\documentclass[11pt,a4paper]{article}
\usepackage[T1]{fontenc}
\usepackage[utf8]{inputenc}
\usepackage{lmodern}
\usepackage[a4paper,margin=0.4in]{geometry}
\usepackage{microtype}
\usepackage{enumitem}
\usepackage{amsmath,amssymb}
\usepackage{mathtools}
\usepackage{hyperref}
\hypersetup{
	colorlinks=true,
	linkcolor=black,
	urlcolor=blue,
	citecolor=black
}

\setlist[itemize]{leftmargin=0.3em}
\setlist[enumerate]{leftmargin=0.3em}

\newcommand{\Title}[1]{\section*{#1}\vspace{-0.25em}\hrule\vspace{0.8em}}
\newcommand{\Field}[2]{\textbf{#1:}~#2\par}
\newcommand{\Affil}[1]{\textit{#1}\par}
\newcommand{\Abstract}{\vspace{0.6em}\textbf{Abstract.}~}

\begin{document}
	
	\begin{center}
		{\LARGE Online Algebra, Logic, and Topology Seminar}\\[0.25em]
		{\large Category Theory Seminar — Abstract}\\[0.5em]
		January 13, 2026
	\end{center}
	
	\vspace{1.25em}\hrule\vspace{1.25em}
	
	\Title{On the commuting tensor product of symmetric multicategories and their bimodules}
	
	\Field{Speaker}{Nicola Gambino}
	\Affil{The University of Manchester, UK}
	
	\Field{Date \& Time}{January 13, 2026 (Tuesday) — 3{:}00\,PM (Coimbra time)}
	\Field{Place}{Google Meet}
	
	\Abstract
We provide a unified treatment of several commuting tensor products considered in the literature,
including the tensor product of enriched categories and the Boardman-Vogt tensor product of symmetric multicategories, subsuming work of Elmendorf and Mandell.  

We then show how a commuting tensor product  extends to bimodules, generalising results of Dwyer and Hess. In particular, we construct a double category of symmetric multicategories, symmetric multifunctors and bimodules and show that it admits a symmetric oplax monoidal structure.

These applications are obtained as instances of a general construction of commuting tensor products on double categories of monads, monad morphisms and bimodules.  

	\medskip
	\textit{(Joint work with Richard Garner and Christina Vasilakopoulou)}
	
\end{document}
